\clearpage\thispagestyle{empty}
\pdfbookmark[0]{Quelle zu Pallas, Ceres, Vesta (26. August 1809)}{planets1809aug-sources}
\begin{center}\textbf{Quelle zu Pallas, Ceres, Vesta (26. August 1809)}\end{center}
\noindent Diese KI-Übertragung ins heutige Deutsch folgt der Mitteilung in \href{https://gdz.sub.uni-goettingen.de/id/PPN236007467}{Gauß’ \textit{Werken}, Bd.~VI} (Göttingen, 1874), S.~315. Der Erstabdruck erschien in den \textit{Göttingischen gelehrten Anzeigen}, Stück 136, 26. August 1809, S.~\href{https://api.digitale-sammlungen.de/iiif/image/v2/bsb10538676_00340/full/1600,/0/default.jpg}{1348}, \href{https://api.digitale-sammlungen.de/iiif/image/v2/bsb10538676_00341/full/1600,/0/default.jpg}{1349}, \href{https://api.digitale-sammlungen.de/iiif/image/v2/bsb10538676_00342/full/1600,/0/default.jpg}{1350}. Das digitalisierte Exemplar stammt aus der Bayerischen Staatsbibliothek.
\par\smallskip\noindent Im Erstabdruck stehen die beiden Beobachtungen der Pallas nochmals. Die Werkausgabe verweist dafür auf den Brief vom 14. August. KI-Übertragung und KI-Satz.
\clearpage
